\section{Native approval and binding-scope reproduction}
\label{sec:a11-native-binding}

\textbf{Version and execution scope.}
The corrected experiment is retained in
\nolinkurl{research/ase-x1-x2-2026-10-21/v2/}
and supplied in \texttt{X1-X2-reproduction.zip}. It uses Python 3.11,
LangChain 1.2.10, langchain-core 1.2.13, LangGraph 1.0.8 and
langgraph-checkpoint 4.0.0. All assistant messages, reviewer decisions
and interventions are scripted; no hosted model is called. The protocol
is an amendment made after inspection of the initial experiment. The
earlier version is retained as historical material; the main text uses
the corrected version and its two saved final runs.

\textbf{X1: pending-interrupt routing and explicit edits.}
Sixteen conditions exercise the unchanged official LangChain middleware
and LangGraph interrupt/resume primitives, using author-written graphs,
tools and recorders. Both rounds of each cross-turn test use the same
graph, checkpointer and thread. Records separate native behavior, official
configuration and adapter choices. Each final run contains 18 actual
effects read from SQLite. A new round creates a new pending interrupt:
resuming with the old interrupt identifier produces no new effect and
leaves the current review pending, whereas the matching identifier
allows the current action. Repeating an unkeyed affirmative input at a
new interrupt is a fresh input from the caller. Authorized edits retain
the earlier proposal and approval payload in checkpoint history.
These conditions characterize routing and explicit-edit semantics for
the tested interfaces. They do not exercise same-identifier pending-state
substitution within a shared interrupted thread, Agent Server
authorization policies, or other persistence backends. Matching an
application grant to a candidate during queued handoffs or reuse is a
separate design question.
\texttt{RAW-EVIDENCE.json} preserves the framework identifiers,
requests and checkpoint history; \texttt{OBSERVATIONS.json} records
the relationships compared across runs.

\textbf{X2: binding and execution checks.}
Nine constructed fixtures cross three main binding comparators and one
raw-content diagnostic with two common execution modes: 54 main routes
and 18 diagnostic routes per run. The comparators share a value guard;
the content projection extends the retained context with decision-relevant
fields. Amount normalization follows the experiment's declared two-decimal
rule. These selected comparators do not implement the manuscript's full
$G$, nor do they replace its canonical-JSON equality.
The unchecked executor reads the external document at execution; the
verified executor compares its bytes to the reviewed digest and consumes
those same bytes. Both modes are applied to every comparator.

For context, normalized content, instance and raw-content binding,
respectively, semantic-violation counts are $4,1,1,1$ with unchecked
execution and $3,0,0,0$ with verified execution (nine fixtures per cell).
Corresponding SQLite-write counts are $8,5,2,4$ and $7,4,1,3$, totaling
34 effects per run. The three remaining context violations concern fields
omitted from its retained comparison. The changed-document case is
stopped by the shared executor check in every arm. Instance binding also
rejects three same-semantics submissions with new identifiers, recording
the cost of the selected attribution rule separately from semantic
violations. The scorer reads the saved approved-effect fixtures and
SQLite in read-only mode without importing the gate or its projection.

\textbf{Saved runs and verification.}
The archive's \texttt{v2/runs/final-1/} and
\texttt{v2/runs/final-2/} contain runner metadata, checks, raw evidence,
scores and SQLite databases. Their semantic JSON/CSV and actual SQL rows
agree. Randomized framework identifiers are retained but not compared
byte for byte. The original \texttt{VERIFICATION.json} also records
historical source-preservation checks at experiment completion; those
checks describe that time, before the manuscript was revised.
The archive's \texttt{verify\_saved.py} verifies the preserved artifact
hashes, compares the two saved executions and independently rescores
the X2 databases without asserting that the later manuscript remains
unchanged.

\textbf{Reproduction entry point.}
After extracting the archive, \texttt{REPRODUCE.md} gives commands for
read-only validation and optional fresh execution. The saved evidence can
be checked with \texttt{python verify\_saved.py}; fresh execution uses
the pinned requirements and \texttt{v2/run\_all.py} with new output
directories. New runs do not overwrite the archived results. The
unaltered \texttt{v2/FINAL-MANIFEST.json} pins 111 source and result files
and has SHA-256
\nolinkurl{7e850e40e45f2afeffed1a2ab5eadb62738343043406049f5a35dc09d55cb066}.
These deterministic checks complement the separate paired-agent study;
they measure interface and binding behavior rather than autonomous
recovery or human-review costs.
